% ============================================
% Chapter: System Design
% ============================================

\chapter{System Design}
\label{ch:system-design}

System Design




This chapter translates the functional and non-functional requirements established in
Chapter 3 into a complete, concrete design for GreenLab. The design is organised across
eight sections that follow the system's four-tier structure: hardware layer (Section 4.2),
Python AI service (Section 4.3), ASP.NET Core backend including the database, REST API,
automation engine, and security design (Sections 4.4--4.7), and Angular dashboard
frontend (Section 4.8). Each section specifies the component topology, interface contracts,
data flows, and decision logic with sufficient precision to guide a direct implementation.

Two conventions hold throughout this chapter. First, all four tiers are referred to by their
fixed names --- the hardware layer, the Python AI service, the ASP.NET Core backend,
and the Angular dashboard --- consistent with the terminology established in Chapter 1
and formalised in Chapter 3. Second, every design decision that implements a specific
requirement cites the requirement's identifier (FR-XX-NN or NFR-XX-NN) so that the
design can be traced backward to the requirement it satisfies and forward to the
implementation described in Chapter 5.


\section{Four-Tier Architecture Overview}


GreenLab is structured as a four-tier cyber-physical system. Each tier owns a single, wellbounded responsibility and communicates with its neighbours through a defined
interface, so that a change within one tier --- replacing the YOLO model, redesigning the
dashboard, or extending the database schema --- does not require a corresponding
change in any other tier.

Table 4.1: GreenLab four-tier architecture --- tiers, responsibilities, and technologies.

Tier Name Primary Responsibility Core Technology
Sense the physical laboratory
environment every polling cycle;
ESP32 microcontroller,
actuate relay-controlled appliances
DHT11, HC-SR04, LDR, 2-
Hardware on command received from the
channel relay module,
Layer backend; maintain JWT-
Arduino C++ firmware,
authenticated HTTP communication;
WiFi HTTP client
fall back to local logic on network
loss.




Tier Name Primary Responsibility Core Technology
Receive IP camera stream frames;
run YOLO-based person detection Python 3, FastAPI,
Python AI using ByteTrack persistent tracking; Ultralytics YOLO
Service maintain per-camera state and (yolo11n.pt), OpenCV,
\begin{lstlisting}
return a stable, de-duplicated person   ByteTrack tracker, uvicorn
\end{lstlisting}

count on demand.

Receive sensor telemetry from the
hardware layer; query the AI service
for occupancy count; evaluate the
lighting decision matrix and fan ASP.NET Core 8, Entity
ASP.NET activation table; Framework Core, SQL
Core Server, JWT Bearer
Backend persist all state changes and issue authentication, REST API
authenticated relay commands; (Swagger/OpenAPI)
enforce role-based access control
and maintain an immutable audit
log.

Present real-time laboratory status
to authenticated instructors and
administrators; issue timed manual
Angular 17, TypeScript,
Angular override commands; submit
HTTP Client, reactive forms,
Dashboard laboratory transfer requests; expose
REST API consumer
administrative controls for labs,
devices, users, schedules, and audit
logs.


The communication paths between tiers are: the hardware layer POSTs sensor events to
the backend via HTTP and receives relay instructions in the JSON response body; the
backend GETs the current person count from the Python AI service's /current-count
endpoint on demand; and the Angular dashboard communicates with the backend
exclusively through the authenticated REST API. There is no direct channel between the
hardware layer and the dashboard, or between the hardware layer and the AI service ---
all routing passes through the backend.




\subsection{Data Flow Diagrams}


To complement the tier-level architecture in Table 4.1, this section presents two
complementary data flow diagrams (DFDs) that trace how information moves through
GreenLab at increasing levels of detail. The context-level diagram (Figure 4.1) treats the
system as a single process and defines its boundary with the four external actors, while
the Level-1 diagram (Figure 4.2) decomposes that single process into the five functional
processes and three persistent data stores that realise the tier responsibilities
summarised above.




Figure 4.1: GreenLab context-level data flow diagram (DFD Level 0) --- the system as a
single process bounded by four external actors and three data stores.

The context diagram fixes the system boundary. Two human actors --- the Instructor and
the Administrator --- and two hardware actors --- the ESP32 hardware and the IP camera
--- exchange data exclusively with the GreenLab System process. The instructor issues
manual override commands and laboratory transfer requests and receives lab status and
override confirmations in return; the administrator issues user, device, and schedule
management commands and transfer decisions, and receives lab status, audit log entries,
and alerts. The ESP32 hardware reports sensor events --- motion, distance, LDR,
temperature, and humidity readings --- and receives relay commands for light and fan
control, while the IP camera contributes a continuous video frame stream. Internally, the
system persists data to three stores: the Sensor/Device Data Store, the Schedule \& User
Management DB, and the Audit Log Archive.




Figure 4.2: GreenLab Level-1 data flow diagram --- decomposition of the

GreenLab System processinto five functional processes and three data stores.

The Level-1 diagram decomposes the single context process into the five processes that
map directly onto the tier responsibilities in Table 4.1. Process 1.0 (Authenticate User)
validates instructor and administrator credentials against D2 --- User \& Schedule Data
before routing role-scoped commands onward. Process 2.0 (Process Sensor Event)
receives sensor readings from the ESP32 hardware and the video frame stream from the
IP camera, persists them to D1 --- Sensor/Device Data, and forwards occupancy-relevant
events onward. Process 3.0 (Manage Devices) issues relay commands --- light and fan
control --- back to the ESP32 hardware and handles instructor override requests
validated by Process 1.0. Process 4.0 (Manage Lab Schedule) and Process 5.0 (Generate
Logs \& Alerts) close the loop by managing transfer requests and writing every state
change to D3 --- Audit Log, from which lab status, alerts, and audit records are returned
to both the instructor and the administrator.


\section{Hardware Layer Design}


\subsection{Component Topology}


The hardware layer is built around a single ESP32 microcontroller that reads four
physical sensors and drives two relay channels. The following pin assignments are taken
directly from the implemented firmware and govern both the circuit layout and the
software configuration.


Table 4.2: ESP32 pin assignments and sensor specifications.


Component GPIO Signal Function Threshold / Range
Pin(s) Type
Entry and proximity
TRIG: detection --- triggers a Trigger if distance
HC-SR04 GPIO 4 / Digital
sensor event when an $\le$ 50 cm; timeout at
Ultrasonic ECHO: PWM
object is within 30 000 µs
GPIO 2 threshold distance

Environmental
monitoring --- Temperature 0--50
DHT11 Single-
temperature and $^\circ$C; Humidity 20--
Temperature
\& Humidity
GPIO 5 wire
digital
humidity readings 90% RH; read on
included in every sensor every loop iteration
event payload

Natural-light level
Dark < 1000; Dim
LDR (Light-
GPIO 34 Analogue monitoring ---
1000--2500; Bright >
Dependent
Resistor)
(ADC) 0--4095 determines whether
2500 (12-bit ADC,
artificial lighting is
needed
\section{V reference)}

LOW = ON (activeRelay Channel
GPIO 26
Digital Controls laboratory low relay); HIGH =
1 (Light) output lighting circuit 1 OFF; default HIGH
(safe-off) at startup

LOW = ON (activeRelay Channel Digital Controls laboratory low relay); HIGH =
GPIO 27
2 (Fan) output ventilation fan circuit 1 OFF; default HIGH
(safe-off) at startup

The following figure shows the complete wiring diagram for the hardware layer
prototype.




Figure 4.3: ESP32 hardware layer wiring diagram --- HC-SR04 ultrasonic sensor (top),
DHT11 temperature/humidity sensor (left), LDR light sensor module (right), and two relay
modules (bottom left and right) connected to the central ESP32 microcontroller.

\subsection{Firmware Architecture}


The firmware is structured around three concurrent concerns: JWT lifecycle
management, sensor polling and HTTP event reporting, and a local WebServer that
receives direct relay commands from the backend when the backend needs to override
the automatic state. These three concerns interact through a set of shared state variables
that are read in the main loop and written by either the event-response handler or the
WebServer route handlers.

JWT Authentication Lifecycle

The firmware maintains a JWT token obtained from POST /api/devices/token with a
payload containing labId, deviceId, and a device password. The token is cached in a
global string and its expiry is tracked by recording millis() + 50 minutes at the time of
issue, which approximates the server-side expiry without requiring the firmware to
parse ISO 8601 timestamps. On every outbound HTTP request, ensureTokenValid()
checks whether the cached token has expired; if so, it calls refreshToken() before
proceeding. If a live request receives a 401 Unauthorized response, the firmware
performs a single in-flight token refresh and retries the request with the new token
before treating the failure as terminal. This design satisfies NFR-SEC-01 (authenticated
device communication) without requiring a real-time clock module.

Sensor Polling and Event Reporting

The main loop runs at approximately 100 ms intervals (governed by the terminal
\begin{lstlisting}
delay(100) call). On each iteration, the loop reads the HC-SR04 distance, the LDR analogue
\end{lstlisting}

value, and the DHT11 temperature and humidity. A proximity event is triggered when
distance $\le$ 50 cm AND the previous reading was not near (rising-edge detection via the
wasNear flag) AND at least 2 000 ms have elapsed since the last event (debounce via
COOLDOWN). This edge-detection pattern prevents continuous event flooding while
ensuring the backend is notified of every genuine entry event, satisfying FR-EE-01 and
FR-EE-02.

When a proximity event fires, the firmware calls sendMotionRequest() which POSTs a
JSON payload containing motionDetected, distanceCm, ldrValue, temperatureCelsius, and
humidity to POST /api/devices/{labId}/event. The backend processes this event, evaluates


the automation engine, and returns a JSON response containing: lightOn, fanOn,
relay3On, relay4On, occupancyCount, and motionIgnored. The firmware parses this
response and actuates the relays accordingly. If occupancyCount is 0, both relays are
forced OFF regardless of any other flags, implementing the empty-room safe-off rule.

Local WebServer for Direct Override

In parallel with the main loop, the firmware runs a WebServer on port 80 that exposes
four endpoints: GET /light/on, GET /light/off, GET /fan/on, GET /fan/off. These endpoints
allow the backend to issue direct relay commands to the ESP32 when a manual override
is activated through the dashboard (FR-MO-01). The WebServer is serviced by calling
server.handleClient() on every main loop iteration, so it shares the same thread as the
sensor loop; the design relies on the low duty cycle of sensor reads and the short HTTP
response time of the relay handlers to avoid contention.

Network Resilience and Fallback

If WiFi.status() is not WL_CONNECTED at the top of the main loop, the firmware attempts
reconnection every 5 000 ms and returns early from the loop body without sending any
HTTP requests. If the backend returns an empty response (connection failure, timeout),
the firmware falls back to local control logic: it compares the current LDR reading
against LIGHT_THRESHOLD (2 000) and directly controls Relay 1 based on the result. This
local fallback satisfies NFR-RES-01 (continued basic operation during network loss).


\section{Python AI Service Design}


\subsection{Service Architecture}


The Python AI service is a FastAPI application that runs on port 8000 and exposes two
endpoints: GET /current-count and GET /health. It is designed as a stateful streaming
service: a dedicated per-camera background thread continuously reads frames from a
camera stream, runs YOLO ByteTrack inference, and updates a shared in-memory count
that the endpoint returns synchronously. This design avoids per-request camera
initialisation overhead and provides sub-second response times to the backend's polling
requests.

\subsection{Model and Inference Pipeline}


The service loads a single YOLO model instance (yolo11n.pt --- the nano variant) at
startup using a thread-safe double-checked locking pattern, ensuring the model is loaded
exactly once regardless of concurrent request traffic. The model is placed on the CPU

(MODEL_DEVICE = 'cpu') to maintain compatibility with deployment environments that
do not have a dedicated GPU. Inference runs via model.track() with the following fixed
configuration:

Table 4.3: Python AI service inference configuration parameters.

Parameter Value Rationale
Nano variant --- lowest parameter count; runs
Model yolo11n.pt within CPU budget for real-time building
automation

ByteTrack persistent tracking assigns stable IDs
bytetrack.y across frames, preventing the ID-fragmentation
Tracker
aml overcounting problem documented by Kim et al.
(Section 2.6)

Reduced input size --- halves inference time
imgsz relative to 640 with acceptable accuracy loss for
occupancy counting

Minimum detection confidence; filters spurious
conf (CONFIDENC
\section{low-confidence boxes while retaining seated}

E_THRESHOL D)
occupants partially occluded by desk equipment

Person class only --- eliminates false positives from
classes [0]
chairs, bags, and other indoor objects

Inference runs every 3rd frame, reducing CPU load
FRAME_SKIP by ~67% while maintaining adequate count update
frequency

Enables ByteTrack cross-frame ID persistence ---
persist true
the key to stable, non-inflated counts


\subsection{Camera Source Abstraction}


The service decouples camera URL resolution from the inference loop through the
camera_sources module. The resolve_source() function accepts any of the following
source formats and normalises them to a value that cv2.VideoCapture accepts:


 USB device: usb:0, usb:1 → integer device index

 Network stream (passed through unchanged): http://, https://, rtsp://

 Video file by extension: .mp4, .avi, .mov, .mkv, .flv, .ts

 Bare IPv4 address (DroidCam backward compatibility): 192.168.x.x →
http://192.168.x.x:4747/video

This abstraction means the backend can register any IP camera by simply passing its IP
address to the /current-count?ip= endpoint, without the service requiring per-camera
configuration at deployment time.

\subsection{Camera Lifecycle Management}


Each distinct stream URL is registered in an in-memory _cameras dictionary keyed by
the resolved URL string. A per-camera CameraState dataclass holds the current count,
connection status, a threading.Event that signals when the first frame has been
processed, and a last_access timestamp. A background cleanup worker runs every 60
seconds and removes camera entries that have not been accessed for more than 300
seconds, releasing the camera thread and freeing the VideoCapture resource. This design
ensures the service does not accumulate idle threads for cameras that the backend has
stopped querying.

\subsection{Backend Integration Contract}


The backend integrates with the AI service through a single polling call:

GET http://{AI_SERVICE_HOST}:8000/current-count?ip={CAMERA_IP}\&port={PORT}

The response JSON is:

\begin{lstlisting}
{ "camera": "<resolved_url>", "people_count":   <int>, "connected": <bool> }
\end{lstlisting}

The backend maps the people_count integer to an occupancy classification (Empty = 0,
Low = 1--5, Moderate = 6--15, High > 15) and passes this classification into the lighting
decision matrix and fan activation table described in Section 4.5.


\section{Database Design}


\subsection{Schema Overview}


The GreenLab database contains eleven tables managed by Entity Framework Core with
a SQL Server provider. The schema is structured around three functional domains:


laboratory management (Labs, LabSensors, LabDevices, StaffLabAssignment),
operational control (LabControlOverrides, ManualControlCommands,
LabSectionSchedules, LabChangeRequests), and audit and identity (Users, AuditLogs,
ScheduleImageImports). The _EFMigrationsHistory table is managed automatically by EF
Core and is not part of the application domain.

The following figure shows the complete entity-relationship diagram for the GreenLab
database.




Figure 4.4: GreenLab database entity-relationship diagram --- 11 application tables
spanning laboratory management, operational control, and audit/identity domains.

\subsection{Table Specifications}


Table 4.4: GreenLab database table definitions and key columns.

Table Purpose Key Columns Notable Constraints
Master record for each Id (PK), Name, LightThreshold and
Labs Directory, TempFanOnThreshol
physical laboratory
LightThreshold, d store per-lab


Table Purpose Key Columns Notable Constraints
TempFanOnThresho calibration values;
ld, ClubCode, LabOccupancyCount
CameraResult, is updated on every
LabOccupancyCount sensor event

Role is an enum
Authenticated system Id (PK), Email, (Admin, Staff,
users FullName, Device);
Users
(instructors,administrat PasswordHash, Role, PasswordHash
ors, service accounts) CreatedAt storesBCrypt-hashed
credential; Email is
unique

SensorType
constrains the
Id (PK), LabDeviceId
Registered sensor reading
(FK), LabId (FK),
LabSensors hardware within each interpretation;
Name, SensorType,
laboratory IsActive allows soft-
Unit, IsActive
disable without
deletion

DevicePasswordHas
Id (PK), LabId (FK),
ESP32 or other IoT h enables the device-
Name,
controller nodes level JWT token
LabDevices DevicePasswordHas
registered per endpoint; LastSeen
h, DeviceType,
laboratory updated on every
IsActive, LastSeen
event receipt

ValidFrom/ValidUntil
Id (PK), UserId (FK),
support time-
LabId (FK),
bounded
Tracks which instructor DayOfWeek,
StaffLabAssig assignments;
is currently assigned to StartTime, EndTime,
n ment AssignmentSource
which laboratory ValidFrom,
distinguishes manual
ValidUntil,
vs schedule-import
AssignmentSource
assignments



Table Purpose Key Columns Notable Constraints
LabId (FK),
MotionSensorDisabl
One row per lab;
ed, ForceLightOn,
SetByUserId
Records the current ForceRelayOn,
LabControlOv provides the actor
manual override state ForceRelay2On,
er rides link to AuditLogs;
for each laboratory ForceRelay3On,
UpdatedAt enables
ForceRelay4On,
expiry evaluation
SetByUserId,
UpdatedAt

Target identifies
Id (PK), LabId (FK), which relay;
RequestedByUserId CommandStatus
Log of every manual (FK), Target, tracks
ManualContro
relay command issued RequestedValue, pending/executed/fai
lC ommands
through the dashboard CommandStatus, le d lifecycle;
RequestedAt, AcknowledgedAt
AcknowledgedAt records ESP32
confirmation

Imported from
Id (PK), LabId (FK), ScheduleImageImpo
Timetable sessions
DayOfWeek, rts pipeline; used by
LabSectionSch imported or manually
StartTime, EndTime, the backend to
e dules created for each
CourseName, determine
laboratory
AcademicTerm scheduled-active
periods

Id (PK), Status lifecycle:
RequestedByUserId Pending →
Instructor requests to (FK), FromLabId Approved/Rejected;
LabChangeRe (FK), ToLabId (FK),
transfer a laboratory ReviewedByUserId
q uests SectionStart,
assignment records the
SectionEnd, administrator who
RequestedAt, Status, actioned it
ReviewedByUserId,


Table Purpose Key Columns Notable Constraints
ReviewedAt

Id (PK), FileName, Supports the
StoredPath, schedule-import
Metadata for schedule
ScheduleImag ContentType, workflow (POST
images uploaded for AI-
eI mports FileSizeBytes, /api/admin/schedules
assisted parsing
UploadedByUserId, /i mport-image →
UploadedAt, Notes analyze → apply)

INSERT-only by
design; no UPDATE
Id (PK), Action,
Immutable record of or DELETE
Details, Entity,
AuditLogs every state-changing permissions granted
UserId (FK),
action in the system to the application
CreatedAt
user; CreatedAt is
server- assigned




\section{Automation Engine Design}


\subsection{Occupancy Classification}


The backend maps the raw people_count integer returned by the AI service to a fourlevel occupancy classification before applying either control matrix. This classification
step decouples the relay-control logic from the precise count value, making the decision
rules stable under small frame-to-frame count fluctuations.

Table 4.5: Occupancy classification thresholds.

people_count Relay Behaviour Default
Classification Range
All relays forced OFF --- hardcoded safe-off rule
Empty that cannot be overridden by lighting/fan
matrices

Low 1--5 Lighting decision matrix evaluated with LDR; Fan


people_count Relay Behaviour Default
Classification Range
activation table evaluated with temperature

As Low --- same matrix evaluation; fan may
Moderate 6--15
activate at a lower temperature threshold

As Low/Moderate --- fan activation priority
High > 15 elevated; lighting assumed required regardless of
LDR


\subsection{Lighting Decision Matrix}


The lighting decision matrix combines three inputs --- LDR zone, occupancy
classification, and the current manual override state --- to determine the value of lightOn
in the event response. LDR zones are derived from the per-lab LightThreshold column in
the Labs table, which allows each laboratory to be calibrated independently.

Table 4.6: Lighting decision matrix (LDR Zone × Occupancy Classification → lightOn).

LDR Threshold Empty Low Moderate High
Zone
ADC < OFF
Dark LightThreshold × (safe- ON ON ON
\section{off)}


LightThreshold × OFF
Dim 0.5 -- (safe- ON ON ON
LightThreshold off)

ON (high
OFF OFF (natural OFF (natural
Adequat LightThreshold -- density
(safe- light light
e LightThreshold overrides
off) sufficient) sufficient)
natural light)




× 1.5
Bright > LightThreshold × 1.5 OFF (safe-off) OFF OFF OFF


Manual override (LabControlOverrides.ForceLightOn = true) short-circuits the matrix
entirely and sets lightOn = true regardless of LDR or occupancy. The backend writes an
AuditLog entry for every override activation and deactivation, recording the actor, the
previous state, and the timestamp.

\subsection{Fan Activation Table}


The fan activation table maps temperature zone (derived from the per-lab
TempFanOnThreshold column) and occupancy classification to the value of fanOn in the
event response.

Table 4.7: Fan activation table (Temperature Zone × Occupancy Classification → fanOn).

Temperature Range Empty Low Moderate High
Zone
Cool < TempFanOnThr eshold $-$ 5$^\circ$C OFF OFF OFF OFF

TempFanOnThr eshold $-$ 5$^\circ$C to
Comfortable OFF OFF OFF ON
TempFanOnThr eshold

TempFanOnThr eshold to
Warm OFF ON ON ON
TempFanOnThr eshold + 5$^\circ$C

Hot > TempFanOnThr eshold + 5$^\circ$C OFF ON ON ON


As with lighting, ForceRelay2On = true in LabControlOverrides overrides the fan table
and forces fanOn = true. The Empty safe-off rule applies to fans as it does to lights:
occupancyCount = 0 forces all relays OFF before the fan table is evaluated.


\section{REST API Design}


\subsection{API Structure}


The GreenLab REST API is implemented in ASP.NET Core 8 and documented through
Swagger/OpenAPI. Endpoints are organised into five groups, each with its own controller



and role-based authorization policy: Auth, Devices, Schedules (admin), Staff, and
AdminLabs. All endpoints require HTTPS in production. The lock icon visible in the API

documentation confirms that every endpoint requires a valid JWT Bearer token; the sole
exception is POST /api/Auth/login, which issues the token.

\subsection{Complete Endpoint Catalogue}


Table 4.8: Full GreenLab REST API endpoint catalogue.


Group Meth
od Path Authorised
Role(s) Description
Authenticate a user
with email and
POS None
Auth /api/Auth/login password; returns
T (public)
JWT Bearer token
and expiry

Receive sensor
event from ESP32;
evaluate
automation engine;
\begin{lstlisting}
return relay
\end{lstlisting}

POS Device
Devices /api/devices/{labId}/event command response
T (JWT)
(lightOn, fanOn,
relay3On,
relay4On,
occupancyCount,
motionIgnored)

Issue a JWT token
None to an authenticated
POS
Devices /api/devices/token (device ESP32 device using
T
password) labId, deviceId, and
device password

POS Register a new IoT
Devices /api/devices Admin
T device in the



Group Meth
od Path Authorised
Role(s) Description
system

Retrieve current
Admin /
Devices GET /api/devices/light light state for all
Staff
registered devices

Upload a timetable
Schedule POS /api/admin/schedules/import- image for AI-
Admin
s T image assisted schedule
extraction

Manually submit
Schedule POS /api/admin/schedules/import- parsed schedule
Admin
s T data data for a
laboratory


Schedule List all schedule
GET /api/admin/schedules/imports Admin
s import records

Trigger AI analysis
Schedule POS /api/admin/schedules/analyze- of the most
Admin
s T image recently uploaded
schedule image

Trigger AI analysis
Schedule POS /api/admin/schedules/analyze- of a specific
Admin
s T import/{importId} uploaded schedule
image by import ID

Apply the results of
an analysed
Schedule POS /api/admin/schedules/apply- schedule import to
Admin
s T analysis the
LabSectionSchedul
es table

Staff GET /api/staff/me/current-labs Staff Return the list of

Group Meth
od Path Authorised
Role(s) Description
laboratories
currently assigned
to the
authenticated
instructor

Issue a manual
override command
for the specified
POS
Staff /api/staff/labs/{labId}/override Staff laboratory
T
(activates
ForceLightOn or
ForceRelay2On)

Submit a
POS laboratory transfer
Staff /api/staff/lab-change-requests Staff
T request (Pending
status)

List all transfer
requests submitted
/api/staff/lab-change-requests/
Staff GET Staff by the
mine
authenticated
instructor

AdminLa List all registered
GET /api/admin/labs Admin
bs laboratories

AdminLa POS Create a new
/api/admin/labs Admin
bs T laboratory record

AdminLa Retrieve a single
GET /api/admin/labs/{id} Admin
bs laboratory by ID

AdminLa Update laboratory
PUT /api/admin/labs/{id} Admin
bs configuration


Group Meth
od Path Authorised
Role(s) Description
(name, thresholds,
club code)

Delete a laboratory
AdminLa DELE
/api/admin/labs/{id} Admin record (soft or hard
bs TE
delete per policy)

Add a timetable
AdminLa POS
/api/admin/labs/{id}/sections Admin section directly to a
bs T
laboratory

List all devices
AdminLa
GET /api/admin/labs/{labId}/devices Admin registered to a
bs
laboratory

Register a new
AdminLa POS
/api/admin/labs/{labId}/devices Admin device for a
bs T
laboratory

AdminLa /api/admin/labs/{labId}/devices/ Retrieve a specific
GET Admin
bs {deviceId} device record

AdminLa /api/admin/labs/{labId}/devices/ Update device
PUT Admin
bs {deviceId} configuration

AdminLa DELE /api/admin/labs/{labId}/devices/ Remove a device
Admin
bs TE {deviceId} from a laboratory

Regenerate the
AdminLa POS /api/admin/labs/{labId}/devices/ device password /
Admin
bs T {deviceId}/token JWT credential for
a specific device

Force lights ON
AdminLa POS across all
/api/admin/labs/light/on Admin
bs T laboratories under
administrator


Group Meth
od Path Authorised
Role(s) Description
authority

Force lights OFF
AdminLa POS
/api/admin/labs/light/off Admin across all
bs T
laboratories

Force fans ON
AdminLa POS
/api/admin/labs/fan/on Admin across all
bs T
laboratories

Force fans OFF
AdminLa POS
/api/admin/labs/fan/off Admin across all
bs T
laboratories


\subsection{Sensor Event Request and Response Contract}


The most critical endpoint --- POST /api/devices/{labId}/event --- defines the primary

communication contract between the hardware layer and the backend. The request body
is:

POST /api/devices/{labId}/event

Authorization: Bearer <device_JWT>

Content-Type: application/json

\begin{lstlisting}
{

\end{lstlisting}

"motionDetected": true,
"distanceCm": 35,
"ldrValue": 1450,
"temperatureCelsius": 28.5,
"humidity": 62.0
\begin{lstlisting}
}

\end{lstlisting}

The response body is:

\begin{lstlisting}
{



\end{lstlisting}

"lightOn": true,
"fanOn": false,
"relay3On": false,
"relay4On": false,
"occupancyCount": 4,
"motionIgnored": false
\begin{lstlisting}
}

\end{lstlisting}

The motionIgnored field is set to true when the backend's automation engine determines

that the event should not trigger a relay state change --- for example when a

LabControlOverride with MotionSensorDisabled = true is active for this laboratory. The

firmware treats motionIgnored = true as a terminal condition and does not apply the
lightOn

or fanOn values.


\section{Security Design}


\subsection{Authentication Architecture}


GreenLab uses a dual-credential authentication architecture. Human users (instructors
and

administrators) authenticate against the Users table through POST /api/Auth/login using

email and BCrypt-hashed password. IoT devices authenticate through POST

/api/devices/token using a device-specific password stored in

LabDevices.DevicePasswordHash. Both paths return a JWT Bearer token with a

configurable expiry (50 minutes in the current firmware, matching the server-side token

lifetime).

\subsection{Role-Based Access Control}


The Users.Role column carries one of three values that map directly to ASP.NET Core

authorization policies:


Table 4.9: Role definitions and permitted endpoint groups.

Role Holder Permitted Endpoint Groups Key Restrictions
Faculty IT Full system access; can
All --- Auth, Devices,
Admi administrators override any lab,
Schedules, Staff (read-only
n and lab approve/reject transfer
view), AdminLabs
coordinators requests, manage users

Scoped to assigned labs
Instructors only; cannot access
Staff assigned to one or Auth, Staff endpoints only AdminLabs endpoints;
more laboratories override is time-limited
and logged

Cannot access any
POST /api/devices/token,
Devic ESP32 firmware human-facing endpoints;
POST
e (service account) identity is per-device, not
/api/devices/{labId}/event only
per-user


\subsection{Audit Logging Design}


Every state-changing action --- automated relay command, manual override activation,
override cancellation, transfer request submission, administrator approval, user
creation, device registration, schedule import --- results in an INSERT into the AuditLogs
table. The AuditLogs table carries no UPDATE or DELETE grant for the application
database user, making every record effectively immutable from the application layer.
Each record contains: Action (the operation type), Details (a JSON-serialisable description
of the change), Entity (the affected entity type and ID), UserId (the authenticated actor),
and CreatedAt (server-assigned timestamp). This design satisfies NFR-SEC-02 (tamperresistant audit trail) and supports the administrative audit-log review view described in
Section 4.8.


\section{Angular Dashboard Design}


\subsection{View Architecture}


The Angular dashboard is a single-page application structured around a role-gated
routing module. On login, the JWT response's role claim determines which navigation
group is activated: Staff users see only their assigned laboratory views; Admin users see

both the administrative panel and can access all laboratory monitoring views.
Unauthenticated routes redirect to the login page. The HTTP interceptor attaches the
stored JWT as a Bearer token to every outbound API call and redirects to login on 401
responses.

\subsection{View Inventory}


Table 4.10: Angular dashboard views, their authorised roles, and primary data sources.


View Route Role Primary API Key User Actions
(indicative) Endpoints
Logi Publi POST Enter email and password;
/login
n c /api/Auth/login receive JWT; redirect to role-




appropriate home
view
View live sensor
readings
(temperature,
My Labs --- GET humidity, LDR,
Real-Time /staff/labs Staff /api/staff/me/current- distance), relay
Monitor labs states, occupancy
count, and active
overrides for each
assigned laboratory

Toggle light or fan
relay state; set
Manual POST override duration;
/staff/labs/:id/
Override Staff /api/staff/labs/{labId}/o cancel active
override
Panel v erride override; view
current override
holder and expiry


Transfer /staff/change- POST /api/staff/lab- Select FromLab,
Staff
Request requests/new change-requests ToLab, session time


appropriate home
view
window; submit
request for
Form
administrator
review

View status
(Pending/Approved/
My
/staff/change- GET /api/staff/lab- Re jected) and
Transfer Staff
requests change-requests/mine admin notes for all
Requests
submitted transfer
requests

Monitor all
laboratories
simultaneously; see
occupancy, relay
Admin ---
Admi states, sensor
All Labs /admin/labs GET /api/admin/labs
n readings, and active
Dashboard
overrides at a
glance; quick-access
light/fan master
controls

Edit lab name,
LightThreshold,
Admin --- TempFanOnThresho
Lab Admi GET/PUT
/admin/labs/:id ld , ClubCode;
Configurati n /api/admin/labs/{id}
manage lab-specific
on sensor and device
records


Admin --- GET/POST/PUT/ Register, update, or
Device /admin/labs/:id/ DELETE deregister ESP32
Admi
Manageme devices /api/admin/labs/{labId} devices; regenerate
n
nt / devices/... device credentials;
view LastSeen

appropriate home
view
timestamp

Upload timetable
image; trigger AI
Admin --- /admin/ analysis; review
Admi POST/GET
Schedule schedules/ parsed result; apply
n /api/admin/schedules/...
Import import to
LabSectionSchedule
s


Admin --- POST Force all lights or all
/admin/labs/ Admi fans ON/OFF across
Master Lab /api/admin/labs/light/o
controls n the entire
Controls n| off, fan/on|off
laboratory



estate in one action;
each action is audit-
logged
Filter by date range,
actor, entity type, or
Admin --- Derived from
Admi action; view full detail
Audit Log /admin/audit AuditLogs via admin
n of each log entry;
Viewer endpoint
export as CSV for
compliance reporting

Review pending
Admin --- transfer requests;
Transfer /admin/change- Admi Derived from approve or reject with
Request requests n LabChangeRequests notes; view history of
Review all requests by status
and date




\subsection{Real-Time Update Strategy}


The dashboard uses a polling strategy for real-time status updates: the laboratory
monitoring views poll their respective API endpoints every 5 seconds using Angular's
interval-based observable pattern. This interval is configurable and was chosen to
balance freshness (the sensor-event loop fires approximately every 100 ms on the
firmware side, but the backend aggregates state between polls) against unnecessary
HTTP load. A future enhancement identified in Chapter 7 is to replace polling with
SignalR WebSocket push notifications, which would reduce latency for time-sensitive
override feedback.


\section{Chapter Summary}


This chapter presented the complete system design for GreenLab across all four tiers.

Section 4.2 specified the hardware layer's component topology (five GPIO-connected
components on a single ESP32), pin assignments, firmware architecture including JWT
lifecycle management, event-triggered reporting, local WebServer override reception,
and network fallback logic. Section 4.3 specified the Python AI service's streaming
inference pipeline (yolo11n.pt + ByteTrack, FRAME_SKIP=3, conf=0.5, imgsz=320), camera
source abstraction, and the backend integration contract. Section 4.4 described all eleven
database tables with their key columns and design intent. Section 4.5 defined the
automation engine through two explicit control matrices --- the lighting decision matrix
and the fan activation table --- both anchored to per-lab calibration values stored in the
database. Section 4.6 catalogued all 31 REST API endpoints across five groups with their
HTTP methods, paths, authorised roles, and descriptions, and specified the full sensorevent request/response contract. Section 4.7 described the dual JWT authentication
architecture, the three-role RBAC model, and the INSERT-only audit log design. Section
\section{inventoried the twelve Angular dashboard views with their routes, roles, API sources,}

and user actions, and described the polling-based real-time update strategy.

Chapter 5 takes these design specifications and describes their direct implementation:
representative ESP32 firmware excerpts, ASP.NET Core controller and service class
structures, Python FastAPI route handlers, and Angular component architecture. Every
implementation decision in Chapter 5 traces back to a design decision made in this
chapter.




